\ssr{ВВЕДЕНИЕ}

Целью курсовой работы является разработка программного решения для генерации и трёхмерной визуализации городской среды.
\begin{itemize}
    \item Разработать алгоритм генерации транспортной сети, представленной соединёнными отрезками дорог с не менее чем тремя перекрёстками и отсутствием изолированных участков.
    \item Создать объемные модели зданий с поддержкой не менее пяти архитектурных типов, обеспечивающую размещение объектов вдоль дорог без взаимных пересечений и минимальную детализацию двадцать полигонов на модель.
    \item Реализовать механизм воспроизводимой генерации городской среды с фиксацией начального значения генератора псевдослучайных чисел.
    \item Разработать модуль интерактивного управления виртуальной камерой, поддерживающий операции перемещения, вращения и масштабирования сцены в реальном времени.
    \item Интегрировать алгоритм карты теней для расчёта освещения с возможностью настройки параметров источников света.
    \item Провести анализ производительности решения, исследовав зависимость времени отрисовки кадра от разрешения изображения и количества полигонов в экранном пространстве.
\end{itemize}

\clearpage